\documentclass[a4paper,fontset=windows]{ctexart}

\usepackage{anyfontsize}
\usepackage{amsmath,amssymb}
\usepackage{graphicx,caption,subcaption}
\usepackage{bm}
\usepackage{geometry,parskip}
\geometry{top=3cm,bottom=3cm,left=2.75cm,right=2.75cm}
\usepackage[shortlabels]{enumitem}
\usepackage{indentfirst}
\setlength{\parindent}{0em}

\begin{document}

\title{\textbf{实验十三\ \ 弦上驻波实验}}
\author{1900011604 张植竣}
\date{2021年5月24日}

\maketitle

\section*{1\ \ 实验数据和现象}

\subsection*{1.1\ \ 基本参数测量}

本实验直接和间接测量的基本参数有：砝码的质量$m$、样品弦的长度$l_0$、样品弦的直径$d_0$、样品弦的质量$m_0$、实验用弦的平均直径$\overline{d}$和实验用弦的弦密度$\rho$，测量数据如下：

\begin{table}[h!]
    \begin{center}
        \begin{tabular}{|c|c|c|c|c|c|}
            \hline
            $m/\mathrm{kg}$ & $l_0/\mathrm{m}$ & $d_0/\mathrm{mm}$ & $m_0/\mathrm{g}$ & $\overline{d}/\mathrm{mm}$ & $\rho/\mathrm{g\cdot m^{-1}}$ \\
            \hline
            0.99999 & 0.694 & 0.302 & 2.34 & 0.298 & 3.283021345 \\
            \hline
            $\sigma_{m_0}/\mathrm{kg}$ & $\sigma_{l_0}/\mathrm{m}$ & $\sigma_{d_0}/\mathrm{mm}$ & $\sigma_{m_0}/\mathrm{g}$ & $\sigma_{\overline{d}}/\mathrm{mm}$ & $\sigma_\rho/\mathrm{g\cdot m^{-1}}$ \\
            \hline
            0.010408330 & 0.005 & 0.023094011 & 0.010408330 & 0.023094011 & 0.715410271 \\
            \hline
        \end{tabular}
    \end{center}
\end{table}

其中实验用弦的弦密度：
\[
\rho = \frac{m_0}{l_0}\times\left(\frac{\overline{d}}{d_0}\right)
\]
其不确定度为：
\[
\sigma_\rho = \rho\sqrt{\left(\frac{\sigma_{m_0}}{m_0}\right)^2 + \left(\frac{\sigma_{l_0}}{l_0}\right)^2 + \left(\frac{2\sigma_{\overline{d}}}{\overline{d}}\right)^2 + \left(\frac{2\sigma_{d_0}}{d_0}\right)^2}
\]

对不确定度贡献最大的是$d_0$和$\overline{d}$的测量，其相对不确定度比另外两个物理量高一个数量级。

\subsection*{1.2\ \ $f-n$关系}

固定弦长$L=0.600\,\mathrm{m}$，拉力$T=3mg$（其中$g=9.801\,\mathrm{m/s^2}$），记录信号源的输入频率$f_\mathrm{src}$、示波器测量的输入频率$f_\mathrm{CH1}$和弦的振动频率$f=f_\mathrm{CH2}$，数据如下：

\begin{table}[h!]
    \begin{center}
        \begin{tabular}{|c|c|c|c|c|c|}
            \hline
            $n$ & 1 & 2 & 3 & 4 & 5 \\
            \hline
            $f_\mathrm{src}/\mathrm{Hz}$ & 40.4 & 80.9 & 121.4 & 161.8 & 203.3 \\
            \hline
            $f_\mathrm{CH1}/\mathrm{Hz}$ & 40.4 & 80.9 & 121.4 & 161.8 & 203.3 \\
            \hline
            $f/\mathrm{Hz}$ & 80.77 & 162.3 & 242.7 & 323.3 & 404.8 \\
            \hline
        \end{tabular}
    \end{center}
\end{table}

作最小二乘法拟合：$f = kn+b$（若无特殊说明，物理量均使用SI单位，下同），如下图：

\begin{figure}[h!]
    \centering
    \includegraphics[width=0.8\linewidth]{1.png}
    \caption{$f-n$关系图}
\end{figure}

拟合数据如下：
\[
k = 80.906,\quad b = 0.056,\quad r = 0.999996798
\]
实验中估计$f_\mathrm{src}$的不确定度为$0.2\,\mathrm{Hz}$，则$\sigma_f = 0.4\,\mathrm{Hz}$，由：
\[
\sigma_{k,A} = \sqrt{\frac{1/r^2-1}{N-2}},\quad \sigma_{k,B} = \frac{\sigma_f}{\sqrt{\sum\limits_{n=1}^{N}(n_i-\overline{n})^2}}
\]
得：
\[
\sigma_{k,A} = 0.118203215,\quad \sigma_{k,B} = 0.04,\quad \sigma_k = \sqrt{\sigma_{k,A}^2 + \sigma_{k,B}^2} = 0.124787820
\]
于是有基频：
\[
f_0 = (80.91 \pm 0.12)\,\mathrm{Hz}
\]
由于此时$L=\dfrac{\lambda_0}{2}=0.600\,\mathrm{m}$，取$\sigma_L = 1\,\mathrm{mm}$，则波长为：
\[
\lambda_0 = (0.3000\pm0.0005)\,\mathrm{m}
\]
弦线中横波传播的速度为：
\[
v = \lambda_0 f_0 = (24.3\pm0.9)\,\mathrm{m/s}
\]

\paragraph{实验现象}实验中可以看到弦线未达到共振时几乎不振动，从起振到共振的过程中，弦线振动越来越明显，振幅越来越大，达到共振时振幅不随时间衰减，用探测线圈测量振幅或肉眼观察均可看到明显的波腹和波节。

\newpage
\subsection*{1.3\ \ $f-T$关系}

固定弦长$L=0.600\,\mathrm{m}$，谐波次数$n=1$，记录信号源的输入频率$f_\mathrm{src}$、示波器测量的输入频率$f_\mathrm{CH1}$和弦的振动频率$f=f_\mathrm{CH2}$，数据如下：

\begin{table}[h!]
    \begin{center}
        \begin{tabular}{|c|c|c|c|c|c|}
            \hline
            $T/mg$ & 1 & 2 & 3 & 4 & 5 \\
            \hline
            $T/\mathrm{N}$ & 9.800902 & 19.601804 & 29.402706 & 39.203608 & 49.004510 \\
            \hline
            $f_\mathrm{src}/\mathrm{Hz}$ & 24.3 & 33.5 & 40.4 & 46.5 & 51.6 \\
            \hline
            $f_\mathrm{CH1}/\mathrm{Hz}$ & 24.3 & 33.5 & 40.4 & 46.5 & 51.6 \\
            \hline
            $f/\mathrm{Hz}$ & 48.59 & 67.02 & 80.77 & 92.76 & 103.1 \\
            \hline
        \end{tabular}
    \end{center}
\end{table}

作最小二乘法拟合：$\ln{f}=k\ln{T}+b$，如下图：

\begin{figure}[h!]
    \centering
    \includegraphics[width=0.8\linewidth]{2_1.png}
    \caption{$\ln{f}-\ln{T}$关系图}
\end{figure}

拟合数据如下：
\[
k = 0.466982332,\quad b = 2.816135830,\quad r = 0.999972830
\]

可见$k \approx 0.5$，说明大致有$f \propto \sqrt{T}$，作最小二乘法拟合：$f=k\sqrt{T}+b$，如下图：

\begin{figure}[h!]
    \centering
    \includegraphics[width=0.8\linewidth]{2_2.png}
    \caption{$f-\sqrt{T}$关系图}
\end{figure}

拟合数据如下：
\[
k = 14.07773482,\quad b = 4.562251877,\quad r = 0.999989629
\]
相关系数$r$几乎等于1，说明$f \propto \sqrt{T}$，与理论推测$f = \dfrac{n}{2L}\sqrt{\dfrac{T}{\rho}} \propto \sqrt{T}$符合。

实验中估计$\sigma_f = 0.4\,\mathrm{Hz}$，这给出：
\[
\sigma_{k,A} = 0.037017274,\quad \sigma_{k,B} = 0.0430830065,\quad \sigma_k = 0.056801620
\]
于是有弦的线密度：
\[
\rho' = \left(\frac{n}{2Lk}\right) = (3.50\pm0.03)\,\mathrm{g/m}
\]

\subsection*{1.4\ \ $f-L$关系}

固定拉力$T=3mg$，谐波次数$n=1$，记录信号源的输入频率$f_\mathrm{src}$、示波器测量的输入频率$f_\mathrm{CH1}$和弦的振动频率$f=f_\mathrm{CH2}$，数据如下：

\begin{table}[h!]
    \begin{center}
        \begin{tabular}{|c|c|c|c|c|c|c|}
            \hline
            $L/\mathrm{m}$ & 0.600 & 0.500 & 0.429 & 0.375 & 0.333 & 0.300 \\
            \hline
            $f_\mathrm{src}/\mathrm{Hz}$ & 40.4 & 48.8 & 56.5 & 64.4 & 73.2 & 80.9 \\
            \hline
            $f_\mathrm{CH1}/\mathrm{Hz}$ & 40.4 & 48.8 & 56.5 & 64.4 & 73.2 & 80.9 \\
            \hline
            $f/\mathrm{Hz}$ & 80.77 & 97.65 & 113.1 & 129.1 & 146.1 & 161.1 \\
            \hline
        \end{tabular}
    \end{center}
\end{table}

作最小二乘法拟合：$\ln{f}=k\ln{L}+b$，如下图：

\begin{figure}[h!]
    \centering
    \includegraphics[width=0.8\linewidth]{3_1.png}
    \caption{$\ln{f}-\ln{L}$关系图}
\end{figure}

拟合数据如下：
\[
k = -0.986322711,\quad b = 3.895703208,\quad r = -0.999949910
\]

可见$k \approx -1$，说明大致有$f \propto L^{-1}$，作最小二乘法拟合：$f=kL^{-1}+b$，如下图：

\begin{figure}[h!]
    \centering
    \includegraphics[width=0.8\linewidth]{3_2.png}
    \caption{$f-L^{-1}$关系图}
\end{figure}

拟合数据如下：
\[
k = 47.941683560,\quad b = 1.548754959,\quad r = 0.999935222
\]
相关系数$r$几乎等于1，说明$f \propto L^{-1}$，与理论推测$f = \dfrac{n}{2L}\sqrt{\dfrac{T}{\rho}} \propto L^{-1}$符合。

实验中估计$\sigma_f = 0.4\,\mathrm{Hz}$，这给出：
\[
\sigma_{k,A} = 0.272854858,\quad \sigma_{k,B} = 0.205313638,\quad \sigma_k = 0.3414724925
\]

\newpage
\section*{2\ \ 分析与讨论}

\paragraph{弦线共振的判据}移走驱动线圈，用手轻轻拨动弦线，探测线圈探测到弦线的振动，示波器上测量得到的频率即是弦的基频，弦线共振时振动频率应约等于该频率，且此时示波器上测量得到的弦振动频率达到最大值。

\paragraph{主要误差来源}
\begin{itemize}
    \item 没有将张力杠杆调节水平会导致弦线张力偏小（因为砝码拉力的力臂减小），共振频率降低。
    
    \item 测量弦的线密度时，没有将弦线拉直（实际上拉直所需要的力非常大，操作上很难做到）会导致弦线的长度测量值偏小，使线密度测量值偏大。弦的直径较小，其相对误差最大，贡献了线密度不确定度的大部分。
    
    \item 探测线圈和驱动线圈之间有电磁干扰（互感），会影响弦线的振动频率测量。
    
    \item 示波器测量得到的信号发生器频率和弦振动频率均有很大的跳跃，弦振动频率极差可以超过$1\,\mathrm{Hz}$。且共振频率测量的可重复性也有约$0.4\,\mathrm{Hz}$
\end{itemize}

\end{document}